%%
%% XXII_introduction_draft.tex
%%
%% Draft Introduction for Paper XXII:
%% "Edge Coercivity and the Airy Boundary Operator:
%%  A Perturbative Stability Theory for Universal Boundary Operators"
%%
%% This introduction presents the paper's results honestly:
%% -- what is proved (Level 1 + 2, unconditional)
%% -- what is conditional (Level 3, rate refinement)
%% -- what the integrable structure contributes (coordinates, not core)
%% -- the RG interpretation (fixed point + irrelevant perturbation)
%%
\documentclass[12pt,a4paper]{amsart}
\usepackage{amsmath,amssymb,amsthm,mathtools,enumitem,booktabs,hyperref}

\newcommand{\DAi}{D^{\mathrm{Airy}}}
\newcommand{\KAi}{K^{\mathrm{Airy}}}
\newcommand{\PAi}{\Pi_{\mathrm{Airy}}}
\newcommand{\Aarith}{A^{\mathrm{arith}}}
\newcommand{\norm}[1]{\|#1\|}

\title{Edge Coercivity and the Airy Boundary Operator\\
{\normalsize A Perturbative Stability Theory for Universal Boundary Operators}\\
\vspace{6pt}
{\small Paper XXII in the Prolate Gram Coercivity Programme}}
\author{Sebastian Schmalnauer}
\date{2026}

\begin{document}
\maketitle

%%=================================================================
\section*{Introduction}
%%=================================================================

\subsection*{The main result}

Let $G_{N,c}$ denote the Gram matrix of the prolate spheroidal wave
functions $\{\psi_k(\,\cdot\,;c)\}_{k=0}^{N-1}$ of bandwidth $c$,
sampled at the Airy-rescaled primes $\{p_j^{\mathrm{Ai}}\}$ near the
spectral edge (see Section~2 for definitions).
The central question of this paper is:

\begin{quote}
  \textit{Is the smallest eigenvalue $\lambda_{\min}(G_{N,c})$
  bounded below by a positive quantity, uniformly as $N, c \to \infty$
  in the Landau--Widom scaling regime?}
\end{quote}

We answer this question affirmatively and unconditionally:

\begin{theorem*}[Main Theorem, P13]
There exist constants $\alpha > 0$ and $c_0 < \infty$ such that
for all $c \ge c_0$ and $N = \lfloor\alpha c\rfloor$:
\[
  \lambda_{\min}(G_{N,c}) \ge \alpha(c) \cdot c^{-1/2} > 0.
\]
\end{theorem*}

The bound is not sharp in its rate;
under the conditional Conjecture~3.3
(convergence of the RHP solution, Section~3),
the rate improves to $c^{-1/3}$.
But the \emph{existence} of a positive lower bound is unconditional.

\subsection*{What drives the result}

The proof rests on a single structural observation:
the mismatch operator
\[
  \DAi := \PAi(\Aarith - \KAi)\PAi
\]
between arithmetic sampling and the Airy kernel operator decomposes as
\begin{equation}\label{eq:intro_decomp}
  \DAi = \underbrace{\PAi(I - \KAi)\PAi}_{\text{universal part}}
       + \underbrace{\PAi(\Aarith - I)\PAi}_{\text{arithmetic perturbation}}.
\end{equation}
The two terms in \eqref{eq:intro_decomp} live on separated asymptotic scales:
\begin{alignat}{3}
  \delta_{\mathrm{Airy}}
  &:= \inf\sigma\bigl(\PAi(I-\KAi)\PAi\bigr)
  &&\ge F_{\mathrm{TW}}(0) \approx 0.96,
  \qquad &&\text{(universal, $c$-independent)}
  \label{eq:intro_delta}\\
  \varepsilon_c
  &:= \norm{\PAi(\Aarith - I)\PAi}_{\mathrm{op}}
  &&\le C\Bigl(\tfrac{\log\log c}{\log c}\Bigr)^{1/2} \to 0
  \qquad &&\text{(arithmetic, vanishing)}
  \label{eq:intro_eps}
\end{alignat}
and the gap persists by operator stability:
$\inf\sigma(\DAi) \ge \delta_{\mathrm{Airy}} - \varepsilon_c > 0$ for large $c$.

The bound \eqref{eq:intro_delta} is a consequence of
the Tracy--Widom theory of the GUE edge~\cite{TW1994}.
The bound \eqref{eq:intro_eps} follows from the
Montgomery--Vaughan mean square theorem for prime exponential
sums~\cite{MV1974}, applied to the $L^2$-Fourier expansion of
Airy-bandlimited test functions.
No Riemann Hypothesis and no conjecture on primes in short intervals
is required.

\subsection*{Universality}

The scale separation \eqref{eq:intro_delta}--\eqref{eq:intro_eps}
establishes a form of \emph{perturbative universality} for the
Airy boundary operator:

\begin{itemize}
  \item The spectral gap \emph{exists} for any positive-density
    sampling sequence (weak universality).
  \item Its leading value $\delta_{\mathrm{Airy}} \approx 0.96$
    is \emph{independent of the prime distribution}
    (universality of the leading term).
  \item The arithmetic corrections $\varepsilon_c \to 0$:
    the prime structure is \emph{irrelevant} in the asymptotic regime.
\end{itemize}

This is not strong universality in the sense of random matrix theory
(where the limiting distribution is completely distribution-free).
It is a weaker, operator-level statement:
the Airy geometry dominates the arithmetic fluctuation asymptotically,
and the spectral gap of $\DAi$ converges to the universal value
$\delta_{\mathrm{Airy}}$ as $c \to \infty$.

In the language of renormalisation theory:
$\KAi$ is the \emph{stable fixed point} of the programme;
$\Aarith - I$ is an \emph{irrelevant perturbation};
and the main theorem is the statement that the fixed point is stable
under this perturbation for all large $c$.

\subsection*{The role of integrable structure}

Paper~XXII is the final paper in a programme that began with
IIKS representations and Riemann--Hilbert analysis (Papers~XII--XXI).
A central observation of this paper is that the main theorem
\emph{does not require} the integrable machinery in its proof.

The IIKS identity, the RHP for $Y_c$, and the Deift--Zhou
steepest descent all appear in Section~3 (the RHP arm of the argument).
They are used to prove Conjecture~3.3
($Y_c \to Y_{\mathrm{Airy}}$ in the appropriate sense),
which in turn gives the rate improvement $\lambda_{\min}(G_{N,c}) \ge \alpha c^{-1/3}$.
But this rate improvement is conditional, and does not affect the
\emph{existence} of the gap.

The integrable structure played an essential \emph{epistemic} role:
without Papers~XII--XXI, we would not have identified $\DAi$ as the
relevant object, nor known that $F_{\mathrm{TW}}(0)$ is the correct
universal constant.
But once the object is identified, the proof of gap-existence
reduces to two classical inputs from outside the integrable programme:
Tracy--Widom theory and the Montgomery--Vaughan mean square.

This separation --- integrable structure as discovery mechanism,
universal spectral geometry as proof core --- is itself a structural
result of the paper.
It suggests that perturbative stability of universal boundary operators
may be a more general phenomenon, accessible by the methods of
this paper without requiring integrable structure case by case.

\subsection*{Organisation of the paper}

The paper is organised to reflect the three-level architecture:

\begin{description}
  \item[Section~2 (Setup).]
    Operators, Airy scaling, and the decomposition~\eqref{eq:intro_decomp}.
    All definitions are independent of IIKS.

  \item[Section~3 (RHP arm, conditional).]
    The Riemann--Hilbert analysis: Assumptions A and B, Conjecture~3.3,
    and the conditional rate improvement to $c^{-1/3}$.
    This section can be read independently and contributes P10.
    It is not logically required for the Main Theorem.

  \item[Section~4 (Discrete arm, O4).]
    The Jaffard--Gram bound: off-diagonal decay $|G_{mn}| \le C/|m-n|^{3/2}$
    via triple-scale stationary phase.
    Contributes P9, which feeds R1 $\to$ P13 via Paper~13b.

  \item[Section~5 (Spectral arm, O5).]
    The gap-persistence theorem: $\delta_{\mathrm{Airy}} - \varepsilon_c > 0$
    via Tracy--Widom (Level~1) and Montgomery--Vaughan (Level~2).
    Unconditional. Closes P12 and contributes the main input to P13.

  \item[Section~6 (Main Theorem).]
    Combines P9 $\to$ R1 $\to$ P13 (Arm~1) and
    O5 $\to$ P12 $\to$ P13 (Arm~2) into the final argument.

  \item[Appendix (Numerical verification).]
    Explicit computation of $\varepsilon_c$ for $c \in [100, 10^4]$;
    verification that $\mu_A/\mu_K \in [1.88, 4.21]$ across this range;
    computation of $C_{\mathrm{crit}} = F_{\mathrm{TW}}(0)/(\zeta(3/2)-1) \approx 0.595$.
\end{description}

\subsection*{Relation to earlier papers}

The present paper concludes the programme begun in Paper~XII.
The key imports from earlier papers are:
\begin{itemize}
  \item Papers~XVIII--XX: the Airy limit for PSWF near the spectral edge,
    establishing that $K_c \to \KAi$ in the appropriate sense.
    Used to \emph{define} $\DAi$; not used in the gap-persistence proof.
  \item Paper~XXI: the bulk spectral approximation ($\sigma(A_{N,c}) \approx \sigma(K_c)$)
    and the existence of at least one positive eigenvalue of $D_{\mathrm{Edge}}$.
    Used as black boxes BB1, BB2.
  \item Paper~13b: the reduction functor
    $\lambda_{\min}(G_{N,c}) \ge c^{-1/2} \Rightarrow \text{spectral gap} \ge C c^{-1/2}$.
    Used as black box BB3 = R1.
\end{itemize}
All other results of this paper are new.

\subsection*{Notation}
Standard: $\langle \cdot, \cdot \rangle$ for $L^2([0,R])$ inner product;
$\|\cdot\|_{\mathrm{op}}$ for operator norm; $\|\cdot\|_{S_1}$ for trace norm;
$\Pi_{\mathrm{Airy}}$ for the spectral projection onto the Airy edge subspace;
$F_{\mathrm{TW}}$ for the Tracy--Widom GUE distribution.
All other notation is defined at first use.

\begin{thebibliography}{99}
\bibitem{TW1994}
C.~Tracy and H.~Widom,
\textit{Level-spacing distributions and the Airy kernel},
Comm.\ Math.\ Phys.\ \textbf{159} (1994), 151--174.
\bibitem{MV1974}
H.~L.~Montgomery and R.~C.~Vaughan,
\textit{Hilbert's inequality},
J.\ London Math.\ Soc.\ (2) \textbf{8} (1974), 73--82.
\bibitem{Bornemann2010}
F.~Bornemann,
\textit{On the numerical evaluation of Fredholm determinants},
Math.\ Comp.\ \textbf{79} (2010), 871--915.
\bibitem{Slepian1978}
D.~Slepian,
\textit{Prolate spheroidal wave functions, Fourier analysis and uncertainty V},
Bell Syst.\ Tech.\ J.\ \textbf{57} (1978), 1371--1430.
\bibitem{Widom1964}
H.~Widom,
\textit{Asymptotic behavior of the eigenvalues of certain integral equations II},
Arch.\ Rational Mech.\ Anal.\ \textbf{17} (1964), 215--229.
\bibitem{PaperXXI}
S.~Schmalnauer,
\textit{Bulk Spectral Approximation for PSWF Gram Matrices} (Paper~XXI),
preprint, 2025.
\bibitem{Paper13b}
S.~Schmalnauer,
\textit{Gram Coercivity and Spectral Gap: The Reduction Functor} (Paper~13b),
preprint, 2024.
\end{thebibliography}

\end{document}
